% \mag 幾何の検証プローブ（docs/pdftex-compatibility.md 更新用の実測根拠）。
% 使い方:
%   pdflatex mag-geometry-probe.tex        → \mag=1000 の基準値
%   pdflatex '\def\MAG{1200}% \mag 幾何の検証プローブ（docs/pdftex-compatibility.md 更新用の実測根拠）。
% 使い方:
%   pdflatex mag-geometry-probe.tex        → \mag=1000 の基準値
%   pdflatex '\def\MAG{1200}% \mag 幾何の検証プローブ（docs/pdftex-compatibility.md 更新用の実測根拠）。
% 使い方:
%   pdflatex mag-geometry-probe.tex        → \mag=1000 の基準値
%   pdflatex '\def\MAG{1200}% \mag 幾何の検証プローブ（docs/pdftex-compatibility.md 更新用の実測根拠）。
% 使い方:
%   pdflatex mag-geometry-probe.tex        → \mag=1000 の基準値
%   pdflatex '\def\MAG{1200}\input{mag-geometry-probe.tex}' → \mag=1200
%   uplatex + dvipdfmx でも同じ手順で比較可（\typeout行は共通。
%   uplatexでは\pdf*系レジスタが未定義のためガードしてある）。
% 要点: jlreq.cls の d 分岐で \pdfhorigin/\pdfvorigin=1truein と
% \pdfpagewidth/height ← (1000/\mag)×\paperwidth を行うため，
% \mag=1200 でも物理用紙は A4 固定・原点 1in 固定で版面オフセットのみ
% ×1.2 される（uplatex+dvipdfmx と同じ意味論。最初の語の x は表組み
% 座標として uplatex と一致する）。
\ifdefined\MAG\mag=\MAG\relax\fi
\documentclass[a4paper]{jlreq}% エンジン自動判定（pdflatex/uplatex共用）
\begin{document}
\typeout{MAG:\the\mag}
\typeout{PW:\the\paperwidth}
\typeout{PH:\the\paperheight}
\ifdefined\pdfpagewidth\typeout{PPW:\the\pdfpagewidth}\fi
\ifdefined\pdfpageheight\typeout{PPH:\the\pdfpageheight}\fi
\ifdefined\pdfhorigin\typeout{HOR:\the\pdfhorigin}\fi
\ifdefined\pdfvorigin\typeout{VOR:\the\pdfvorigin}\fi
\typeout{OSM:\the\oddsidemargin}
\typeout{ESM:\the\evensidemargin}
\typeout{TMGM:\the\topmargin}
\typeout{HOFF:\the\hoffset}
\typeout{VOFF:\the\voffset}
\typeout{TW:\the\textwidth}
\noindent MARKa 本文テスト
\end{document}
' → \mag=1200
%   uplatex + dvipdfmx でも同じ手順で比較可（\typeout行は共通。
%   uplatexでは\pdf*系レジスタが未定義のためガードしてある）。
% 要点: jlreq.cls の d 分岐で \pdfhorigin/\pdfvorigin=1truein と
% \pdfpagewidth/height ← (1000/\mag)×\paperwidth を行うため，
% \mag=1200 でも物理用紙は A4 固定・原点 1in 固定で版面オフセットのみ
% ×1.2 される（uplatex+dvipdfmx と同じ意味論。最初の語の x は表組み
% 座標として uplatex と一致する）。
\ifdefined\MAG\mag=\MAG\relax\fi
\documentclass[a4paper]{jlreq}% エンジン自動判定（pdflatex/uplatex共用）
\begin{document}
\typeout{MAG:\the\mag}
\typeout{PW:\the\paperwidth}
\typeout{PH:\the\paperheight}
\ifdefined\pdfpagewidth\typeout{PPW:\the\pdfpagewidth}\fi
\ifdefined\pdfpageheight\typeout{PPH:\the\pdfpageheight}\fi
\ifdefined\pdfhorigin\typeout{HOR:\the\pdfhorigin}\fi
\ifdefined\pdfvorigin\typeout{VOR:\the\pdfvorigin}\fi
\typeout{OSM:\the\oddsidemargin}
\typeout{ESM:\the\evensidemargin}
\typeout{TMGM:\the\topmargin}
\typeout{HOFF:\the\hoffset}
\typeout{VOFF:\the\voffset}
\typeout{TW:\the\textwidth}
\noindent MARKa 本文テスト
\end{document}
' → \mag=1200
%   uplatex + dvipdfmx でも同じ手順で比較可（\typeout行は共通。
%   uplatexでは\pdf*系レジスタが未定義のためガードしてある）。
% 要点: jlreq.cls の d 分岐で \pdfhorigin/\pdfvorigin=1truein と
% \pdfpagewidth/height ← (1000/\mag)×\paperwidth を行うため，
% \mag=1200 でも物理用紙は A4 固定・原点 1in 固定で版面オフセットのみ
% ×1.2 される（uplatex+dvipdfmx と同じ意味論。最初の語の x は表組み
% 座標として uplatex と一致する）。
\ifdefined\MAG\mag=\MAG\relax\fi
\documentclass[a4paper]{jlreq}% エンジン自動判定（pdflatex/uplatex共用）
\begin{document}
\typeout{MAG:\the\mag}
\typeout{PW:\the\paperwidth}
\typeout{PH:\the\paperheight}
\ifdefined\pdfpagewidth\typeout{PPW:\the\pdfpagewidth}\fi
\ifdefined\pdfpageheight\typeout{PPH:\the\pdfpageheight}\fi
\ifdefined\pdfhorigin\typeout{HOR:\the\pdfhorigin}\fi
\ifdefined\pdfvorigin\typeout{VOR:\the\pdfvorigin}\fi
\typeout{OSM:\the\oddsidemargin}
\typeout{ESM:\the\evensidemargin}
\typeout{TMGM:\the\topmargin}
\typeout{HOFF:\the\hoffset}
\typeout{VOFF:\the\voffset}
\typeout{TW:\the\textwidth}
\noindent MARKa 本文テスト
\end{document}
' → \mag=1200
%   uplatex + dvipdfmx でも同じ手順で比較可（\typeout行は共通。
%   uplatexでは\pdf*系レジスタが未定義のためガードしてある）。
% 要点: jlreq.cls の d 分岐で \pdfhorigin/\pdfvorigin=1truein と
% \pdfpagewidth/height ← (1000/\mag)×\paperwidth を行うため，
% \mag=1200 でも物理用紙は A4 固定・原点 1in 固定で版面オフセットのみ
% ×1.2 される（uplatex+dvipdfmx と同じ意味論。最初の語の x は表組み
% 座標として uplatex と一致する）。
\ifdefined\MAG\mag=\MAG\relax\fi
\documentclass[a4paper]{jlreq}% エンジン自動判定（pdflatex/uplatex共用）
\begin{document}
\typeout{MAG:\the\mag}
\typeout{PW:\the\paperwidth}
\typeout{PH:\the\paperheight}
\ifdefined\pdfpagewidth\typeout{PPW:\the\pdfpagewidth}\fi
\ifdefined\pdfpageheight\typeout{PPH:\the\pdfpageheight}\fi
\ifdefined\pdfhorigin\typeout{HOR:\the\pdfhorigin}\fi
\ifdefined\pdfvorigin\typeout{VOR:\the\pdfvorigin}\fi
\typeout{OSM:\the\oddsidemargin}
\typeout{ESM:\the\evensidemargin}
\typeout{TMGM:\the\topmargin}
\typeout{HOFF:\the\hoffset}
\typeout{VOFF:\the\voffset}
\typeout{TW:\the\textwidth}
\noindent MARKa 本文テスト
\end{document}
